\documentclass[12pt]{article}

\newif\ifsol
\soltrue

% Unicode compatibility
\usepackage{iftex}
\ifPDFTeX
  \usepackage[utf8]{inputenc}
  \usepackage[noTeX]{mmap}
  \usepackage[T1]{fontenc}
\fi
% XeTeX does not support mmap
\ifLuaTeX
  \usepackage{luatex85}
  \usepackage[noTeX]{mmap}
\fi

\setlength{\textheight}{9in}
\setlength{\textwidth}{7.1in}
\setlength{\evensidemargin}{-0.2in}
\setlength{\oddsidemargin}{-0.2in}
\setlength{\headsep}{30pt}
\setlength{\topmargin}{-0.3in}

\usepackage{amsthm,amsfonts,amsmath,xspace,tikz,varwidth,cancel,soul}
\usepackage{hyperref,verbatim}
\newtheorem{theorem}{Theorem}
\theoremstyle{definition}
\newtheorem{definition}{Definition}

\newcommand{\addmedskip}{\addvspace{\medskipamount}}

\newcommand{\addbigskip}{\addvspace{\bigskipamount}}

\pagestyle{myheadings}
\markboth{BU CAS CS 538.}{BU CAS CS 538.} 
\newcounter{problemnum}
\setcounter{problemnum}{0}
\newenvironment{problem}
     {\addbigskip \setcounter{partnum}{0}
      \noindent\stepcounter{problemnum}\textbf{Problem
                                             \arabic{problemnum}.\ }}
     {\par\addbigskip}
     
\ifsol
\newenvironment{solution}
     {\addbigskip
      \noindent\color{blue}\textbf{Solution.\ }}
     {\par\addbigskip}
\else
\newenvironment{solution}
     {\comment}
     {\endcomment}
\fi

\newcounter{partnum}
\setcounter{partnum}{0}
\newenvironment{ppart}
     {\addmedskip
      \noindent\stepcounter{partnum}\textbf{(\alph{partnum})}\ }
     {\par\addbigskip}

\newcommand{\Enc}{\mathrm{Enc}}
\newcommand{\Dec}{\mathrm{Dec}}
\newcommand{\pk}{\mathit{pk}}
\newcommand{\sk}{\mathit{sk}}
\newcommand{\qrn}{\mathit{QR}_n}
\newcommand{\qrp}{\mathit{QR}_p}
\newcommand{\qrpone}{\mathit{QR}_{p_1}}
\newcommand{\qrptwo}{\mathit{QR}_{p_2}}
\newcommand{\prg}{\mathit{J}}

\newcommand{\crt}{\mathrm{crt}}


\newcommand{\eqdef}{\buildrel{\scriptstyle{\mathit{def}}}\over{=}}
\newcommand{\rfrom}{\buildrel{\scriptstyle{\mathrm{R}}}\over{\leftarrow}}

\newcommand{\mbmod}{\,\%\,}

\newcommand{\A}{\mathcal{A}}
\renewcommand{\L}{\mathcal{L}}
\newcommand{\Z}{\mathbb{Z}}
\newcommand{\G}{\mathbb{G}}
\newcommand{\X}{\mathcal{X}}
\newcommand{\Y}{\mathcal{Y}}
\newcommand{\K}{\mathcal{K}}
\newcommand{\M}{\mathcal{M}}
\newcommand{\E}{\mathcal{E}}
\newcommand{\C}{\mathcal{C}}
\newcommand{\R}{\mathcal{R}}
\renewcommand{\S}{\mathcal{S}}
\newcommand{\I}{\mathcal{I}}
\newcommand{\T}{\mathcal{T}}
\newcommand{\randk}{\mathsf k}
\newcommand{\randm}{\mathsf m}
\newcommand{\randc}{\mathsf c}
\newcommand{\Adv}{\mathcal{A}}
\newcommand{\B}{\mathcal{B}}
\newcommand{\SSadv}{\textrm{SS}\textsf{adv}}
\newcommand{\PRGadv}{\textrm{PRG}\textsf{adv}}
\newcommand{\threebitadv}{\textrm{3Bit}\textsf{adv}}
\newcommand{\PRFadv}{\textrm{PRF}\textsf{adv}}
\newcommand{\MACadv}{\textrm{MAC}\textsf{adv}}
\newcommand{\CPAadv}{\textrm{CPA}\textsf{adv}}
\newcommand{\BCadv}{\textrm{BC}\textsf{adv}}
\newcommand{\CRadv}{\textrm{CR}\textsf{adv}}
\newcommand{\DDHprimeadv}{\textrm{DDH}'\textsf{adv}}
\newcommand{\DDHadv}{\textrm{DDH}\textsf{adv}}
\newcommand{\DDHprime}{DDH$'$}

\newtheorem{lemma}{Lemma}

\newcommand{\indist}{  \mathrel{\vcenter{\offinterlineskip
  \hbox{$\sim$}\vskip-.35ex\hbox{$\sim$}\vskip-.35ex\hbox{$\sim$}}}}

\newcommand{\samples}{\xleftarrow{R}}

\newcommand{\hdl}{H_{\mathrm{dl}}}

%%% BEGIN MIKE ROSULEK'S MACROS
\newcommand{\codebox}[1]{%
        \begin{varwidth}{\linewidth}%
        \begin{tabbing}%
            ~~~\=\quad\=\quad\=\quad\=\kill
            #1
        \end{tabbing}%
        \end{varwidth}%
}

\newcommand{\fcodebox}[1]{%
    \fcolorbox{black}{white}{\codebox{#1}}%
}

\newcommand{\titlecodebox}[2]{%
    \fboxsep=0pt%
    \fcolorbox{black}{black!10}{%
        \begin{varwidth}{\linewidth}%
        \centering%
        \fboxsep=3pt%
        \colorbox{black!10}{#1} \\
        \colorbox{white}{\codebox{#2}}%
        \end{varwidth}%
    }
}



\newcommand{\link}{\diamond}
\newcommand{\lib}[1]{\ensuremath{\L_{\textsf{#1}}}\xspace}
\newcommand{\subname}[1]{\ensuremath{\textsc{#1}}\xspace}

%%% END OF MIKE ROSULEK'S MACROS


\begin{document}
\begin{center}
\Large{\textbf{CAS CS 538.  \ifsol Solutions to \fi Problem Set 9}}\\
\smallskip
\large{\textbf{Due electronically via Gradescope, Monday April 13, 2026 11:59pm
}}
\end{center}

\noindent
\textbf{How to submit:}  See instructions on Problem Set 1. 
 
\section*{An Efficient DL-based PRF}

\textbf{Do not be intimidated by the long problem. It's broken down into four small parts and each part is quite doable. At the end of the four parts, you get a pretty cool theorem.}


The goal of the problem below is to show the following. Let $\G$ be a group of prime order $q$ generated by an element $g$. Let $k$ be some constant. Let $\alpha_0, \dots, \alpha_\ell$ be elements of $\G$. For an input $x$ that consists of $n$ bits $b_1, \dots, b_\ell$, let $F((\alpha_0, \dots, \alpha_\ell), x)$ be computed as follows:
\begin{enumerate}
\item Compute \[\alpha_x=\alpha_0 \cdot \left(\prod_{i \mbox{ such that }b_i = 1} \alpha_i \right)= \alpha_0 \cdot\left( \prod_{i =1}^\ell \alpha_i^{b_i}\right)\]
\item Output $g^\alpha_x$
\end{enumerate}

You will prove that $F((\alpha_0, \dots, \alpha_\ell), x)$ is a PRF, where $(\alpha_0, \dots, \alpha_\ell)$ is the key and $x$ is the input.

The big idea of the proof is the same as the Goldreich-Goldwasser-Micali binary tree construction of Section 4.6: this PRF can be viewed as a binary tree of depth $\ell$, and each level is a PRG. Here is how. Let $h=g^\alpha_0$.

\begin{itemize}
\item The root level contains one node $h$. 
\item The next level contains two nodes $h$ and $h^{\alpha_1}$. 
\item The next level contains four nodes $h, h^{\alpha_2}, h^{\alpha_1}, h^{\alpha_1\alpha_2}$
\item The next level contains eight nodes $h, h^{\alpha_3}, h^{\alpha_2}, h^{\alpha_2\alpha_3}, h^{ \alpha_1},h^{\alpha_1\alpha_3}, h^{\alpha_1\alpha_2},h^{\alpha_1\alpha_2\alpha_3}$.
\item \dots
\item The level at depth $d$ contains $g^{\alpha_x}$ for every binary string $x$ where the first $d$ bits are arbitrary and the last $\ell-d$ bits are 0.
\item\dots
\item The last level contains $g^{\alpha_x}$ for every for every binary string of length $\ell$
\end{itemize}

All we need to prove is that each level $d$ is a secure PRG when the corresponding $\alpha_d$ is kept secret. Thus, let $\prg:\G^k\times\Z_q\rightarrow \G^{2k}$ be defined as follows: on input $(X_1, \dots, X_k, \alpha)$, compute $Y_i=X_i^\alpha$ for $i=1\dots k$. Output $(X_1, Y_1, X_2, Y_2, \dots, X_k, Y_k)$.

\begin{theorem} Under the DDH assumption (Definition 10.8 in Section 10.5) on $\G$, $\prg$ is a secure pseudorandom generator.
\end{theorem}

You will now prove this theorem by proving the following parts.

\begin{problem}  (60 points at 15 each)

\begin{ppart}
Consider the following assumption we call \DDHprime. It is exactly the same as DDH described in Attack Game 10.6 and Definition 10.8, except in Experiment 1, $\gamma$ is chosen at random from $\Z_q-\{\alpha \beta\}$ (call this modified experiment $1'$ and the corresponding advantage $\DDHprimeadv$). Show that assuming DDH holds for $\G$, then so does \DDHprime. (\emph{Hint:} apply the difference lemma to prove that experiments 1 and 1' are indistinguishable.)
\end{ppart}
\begin{solution}
Your solution goes here
\end{solution}

\begin{ppart}
Prove the following.
\begin{lemma}
Fix any $\alpha$, $\beta$, and $\gamma \not\equiv_q \alpha\beta$ in $\Z_q$. Define the map $M: \Z_q\times \Z_q \to \Z_q\times \Z_q$ as follows. Given inputs $d, e$:
\begin{align*}
x&=(d+\beta e)\bmod q\,;\\
y&=(\alpha d+\gamma e)\bmod q\,;\\
&\mbox{Output } (x, y)\,. 
\end{align*}
This map is a bijection.
\end{lemma}

\emph{Hint:} Find the inverse of $M$ by solving a system of two linear equations. Linear algebra extends to modular arithmetic modulo a prime $q$. Note that $\alpha, \beta$, and $\gamma$ are not inputs to the map; rather, they are used to specify a particular map.

\end{ppart}
\begin{solution}
Your solution goes here
\end{solution}

\begin{ppart} 
Prove the following. (\emph{Hint:} work in the exponent space---i.e., consider logarithms of $X$ and $Y$ base $g$ --- and use the previous part. The variable naming here deliberately works well with the previous part.)
\begin{lemma}
For any $u=g^\alpha$, $v=g^\beta$, and $w_1 \ne g^{\alpha\beta}$, if $d$ and $e$ are a uniformly random pair of elements from $\Z_q$, then $X=g^d v^e$ and $Y=u^d w_1^e$ is a uniformly random pair of elements from $\G$ (i.e., every pair $X, Y$ occurs with probability $1/q^2$).
\end{lemma}
\end{ppart}
\begin{solution}
Your solution goes here
\end{solution}

\begin{ppart} 
Use the previous part to build a reduction that shows \DDHprime{} implies the security of $\prg$. (Combined with part (a), this proves Theorem 1.)
In particular, show that for every efficient adversary $\Adv$ breaking $\prg$ with advantage $\epsilon$, there is an efficient adversary $\Adv'$ breaking \DDHprime{} with advantage $\epsilon$ as well.
(\emph{Hint:} This is a single reduction; no hybrid arguments are needed. The previous part tells you exactly how your reduction should transform $u, v$, and $w$ that it receives from the \DDHprime{} challenger. It should that transformation $k$ times.)
\end{ppart}
\begin{solution}
Your solution goes here
\end{solution}


\end{problem}



\section*{Backdoored PRGs}



\begin{problem} (40 points at 20 points each)

    In 2006, the National Institute for Standards and Technology (NIST) unveiled a new standard called \href{https://en.wikipedia.org/wiki/Dual_EC_DRBG}{Dual\_EC\_DRBG}, a PRG construction designed by the NSA. The construction was peculiar. It was clunkier than other existing PRGs, and up to 1000 times slower. Then, in 2007, two Microsoft cryptographers discovered something very unusual: the scheme had the possibility to contain a backdoor, which would allow an adversary to predict all future randomness given some output of the PRG. The existence of such a backdoor was confirmed by the \href{https://en.wikipedia.org/wiki/2010s_global_surveillance_disclosures}{Snowden disclosures} in 2013: the NSA had indeed inserted a backdoor into this PRG.

Dual\_EC\_DRBG is designed as a stateful PRG. At its core is a PRG $G$ which is repeatedly composed with itself sequentially, as in section 3.4.2 of the textbook. We know that this construction is secure as long as the underlying $G$ is a secure PRG.

The following is a simplified version of the Dual\_EC\_DRBG PRG\footnote{The real scheme uses an elliptic curve group instead of $QR_p$, and a different bijection is used. It also removes some bits of the PRG output. What this problem calls $g$ and $h$ are referred to as $Q$ and $P$.}. Let $p$ be a Sophie Germain prime, and $q=(p-1)/2$. Let $g$ and $h$ be generators of $QR_p$, sampled independently and uniformly. Let $\mathrm{abs}: QR_p \to \mathbb{Z}_q$ be the bijection defined in 8.5.1, such that $\mathrm{abs}(z)$ returns $z$ if $z<p/2$ and $-z$ otherwise.
The underlying PRG is as follows:

\begin{figure}[h!]
    
\centering 
\fcodebox{$G(s_0):$\\
\> $r \gets h^{s_0}$\\
\> $s_1 \gets \mathrm{abs}\left( g^{s_0} \right) $ \\
\> return $(r,s_{1})$
}
\end{figure}

\begin{ppart}
    Prove that when $g$ and $h$ are sampled uniformly at random, $G$ is a secure PRG, assuming DDH is hard in $QR_p$. To capture the notion of $g$ and $h$ being generated uniformly at random, we need to add a step to the PRG security game (Attack Game 3.1 in the textbook): we will have the challenger sample $g$ and $h$ and send them to the adversary as the first step, and then use those values when computing $G$ in Experiment 0.

\end{ppart}

\begin{solution}
Your solution goes here
\end{solution}

\begin{ppart}
    Now show that if $g$ and $h$ are generated maliciously, the scheme has a backdoor. In particular, describe how to generate $g$ and $h$ in a way that produces a trapdoor key $sk$ and an efficient function $T$ that takes in $sk$ and $r_i$ (the $i$-th output of $G$ using $g$ and $h$, using an unknown random seed), and computes the next state $s_{i+1}$. (Note that once you know the state, you can compute all future outputs of the PRG.)

    Your maliciously generated $g$ and $h$ should be indistinguishable from an honestly generated $g$ and $h$. This ensures that the existence of an intentional backdoor is undetectable.

\end{ppart}

\begin{solution}
Your solution goes here
\end{solution}



\end{problem}


\end{document}